\section{Method}
\label{sec:method}

Our framework has four components:
(1)~a projection of the observed features onto the Standard Fairness Model
(SFM), which fixes the causal assumptions and the classifier-scale estimands
they identify (Section~\ref{sec:projection});
(2)~a block-ordered generative model of the mediators, $P(W\mid X,Z)$
(Section~\ref{sec:flows});
(3)~an individualized interventional distribution $Q_a$ that simulates the
consequences of an action plan $a$ (Section~\ref{sec:intervention}); and
(4)~a constrained recourse program that returns the cheapest plan whose
prediction distribution is both valid and close, in Wasserstein distance, to
the prediction distribution the recipient would receive under the advantaged
group's mediators. If no plan satisfies both constraints, the program
abstains (Section~\ref{sec:recourse_selection}).
Section~\ref{sec:theory} characterizes exactly when the two constraints can
hold simultaneously. Section~\ref{sec:evaluation} defines how we measure the
disparity before and after recourse.

% ---------------------------------------------------------------------------
\subsection{SFM Projection and Identification}
\label{sec:projection}

Prior causal recourse methods assume a fully specified structural causal model
over all features~\citep{karimi2021algorithmic}. We instead commit only to a
partition of the features into causal roles and an order over groups of
mediators. Leaving the edges within each role unspecified makes this strictly
weaker than a full SCM, although it remains an assumption. We state it
explicitly and separate the parts that identify the disparity from the parts
that are needed only to simulate recourse.

\begin{definition}[SFM projection]
\label{def:sfm}
Partition the observed variables into
\begin{itemize}[nosep]
  \item $X\in\{x^{(0)},x^{(1)}\}$: a binary sensitive attribute (e.g., sex),
        with $x^{(0)}$ the disadvantaged group;
  \item $Z$: pre-treatment covariates that are not descendants of $X$
        (e.g., age, place of birth);
  \item $W$: mediators, i.e., descendants of $X$ that precede the decision
        (e.g., education, occupation, hours worked);
  \item $Y$: the binary outcome, used only to train the classifier and to
        identify recourse recipients.
\end{itemize}
A fixed classifier $\hat f:\mathcal X\times\mathcal W\times\mathcal Z\to[0,1]$
induces the prediction $\hat Y=\hat f(X,W,Z)$, a deterministic child of
$(X,W,Z)$, and a positive decision when $\hat f\ge\tau$. The projected graph
$\mathcal G_{\mathrm{SFM}}$ is complete across clusters. It contains the direct
path $X\to\hat Y$, the indirect paths $X\to W\to\hat Y$ (including chains within
$W$), and the spurious paths $X\leftrightarrow Z\to\hat Y$ and
$X\leftrightarrow Z\to W\to\hat Y$.
\end{definition}

\noindent For $b\in\{x^{(0)},x^{(1)}\}$, let $W_b$ denote the potential
mediator, the value $W$ would take had $X$ been set to $b$. We write
\begin{equation}
  \mu_{ab}(z)\;:=\;\mathbb E\bigl[\hat f(a,W_b,z)\mid Z=z\bigr],
  \qquad a,b\in\{x^{(0)},x^{(1)}\},
  \label{eq:mu_ab}
\end{equation}
where the first index is the value of $X$ entering $\hat f$ directly and the
second is the value of $X$ that generates the mediators.

\begin{assumption}[Identification of classifier-scale effects]
\label{ass:ident}
(i)~$Z$ contains no descendant of $X$;
(ii)~consistency: $W=W_b$ whenever $X=b$;
(iii)~positivity: $P(X=b\mid z)>0$ for all $z$ and $b$;
(iv)~no unmeasured $X$--$W$ confounding given $Z$: $W_b\perp X\mid Z$.
\end{assumption}

\noindent Under Assumption~\ref{ass:ident},
$P(W_b\mid z)=P(W\mid X=b,z)$, so every $\mu_{ab}(z)$ is identified from the
observational mediator distribution:
\begin{equation}
  \mu_{ab}(z)\;=\;\mathbb E_{W\sim P(W\mid X=b,\,z)}\bigl[\hat f(a,W,z)\bigr].
  \label{eq:mu_id}
\end{equation}
Because $\hat f$ is a known function, the step $W\to\hat Y$ cannot be
confounded. No assumption about mediator--outcome confounding or cross-world
independence for $Y$ is needed, so the SFM's restrictions on latent $W$--$Y$
and $X$--$Y$ confounding go unused. This is a deliberate choice of target: we
study the disparity produced by the decision rule, which is what recourse acts
on. Reading these quantities as effects on the true outcome $Y$ would also
require the standard sequential-ignorability conditions, which we do not
claim. Every estimand holds $z$ fixed, so the spurious $X$--$Z$ association is
excluded from the mediated disparity rather than attributed to it. The
mediated disparity is the effect of $X$ on $\hat f$ through the mediator
\emph{set} $W$, and it includes chains within $W$. Effects through a single
mediator would need further path-specific assumptions.

\paragraph{Permitted versus assumed edges.}
The projection allows every mediator to depend on $X$, $Z$, and earlier
mediators, and allows $\hat f$ to depend on all of them. Allowing an edge is
not a substantive assumption. If $X$ has no effect on some mediator block, the
correctly specified conditional does not depend on $x$, and that block
contributes nothing to the mediated disparity. A superfluous edge costs
statistical efficiency, not validity. The causal content lies in the
\emph{omitted} edges:
$X$ and $W$ do not cause $Z$;
there is no unmeasured $X$--$W$ confounding given $Z$;
and, where stated, dataset-specific exclusion restrictions apply (e.g.,
no direct edge from sex to the repayment terms (credit amount, duration) in German Credit;
Appendix~\ref{app:roles}).
The order among mediator blocks, which only recourse requires, is stated
separately in Assumption~\ref{ass:intervene}.

\paragraph{Role assignment.}
The partition is fixed before any model is fit, using a stated protocol rather
than case-by-case judgement.
$X$ is determined by the fairness question.
A variable is placed in $Z$ only if it is determined before $X$ could affect
it, by temporal or biological precedence (e.g., age, place of birth).
A variable is placed in $W$ if the domain literature identifies it as a
consequence of $X$ that precedes the decision. Variables that fit neither role
are excluded, and the reason is recorded (e.g., commute time in ACS).
Which paths count as unfair is a normative choice and is stated per dataset.
Appendix~\ref{app:roles} lists the role of every variable with a one-line
justification. The most damaging misassignment places a pre-treatment
variable that is confounded with $X$ into $W$: its spurious association with
$X$ would then be counted as mediated disparity. The precedence criterion for
$Z$ guards against exactly this error. A variable unrelated to $X$ placed in
$W$ is harmless, since no mediated effect flows through it. Finally, the
semi-synthetic German Credit benchmark satisfies Assumptions~\ref{ass:ident}
and~\ref{ass:intervene} by construction and has known oracle effects, which
lets us check the pipeline when the projection is correct.

Actionability, i.e., which coordinates of $W$ an individual can change and in
which direction, is a property of the recourse problem, not of the causal
model. We specify it in Section~\ref{sec:intervention}.

% ---------------------------------------------------------------------------
\subsection{Normalizing Flow Estimation of $P(W\mid X,Z)$}
\label{sec:flows}

The nested counterfactual predictions in~\eqref{eq:mu_id} and the recourse
program (Section~\ref{sec:recourse_selection}) both require sampling from
$P(W\mid X,Z)$, and recourse additionally requires simulating interventions on
$W$. Mediators may be continuous, ordinal, or categorical, and the SFM itself
imposes no order among them, treating $W$ as an unordered set. Some mediators
have a defensible causal order (education precedes occupation), whereas others
are determined jointly and have none (weekly hours and weeks worked; LSAT and
GPA). We therefore partition $W$ into type-homogeneous \emph{blocks}
$W=(W_1,\ldots,W_L)$, ordered between blocks but unordered within them, and
factorize
\begin{equation}
  P(W\mid X,Z)\;=\;\prod_{\ell=1}^{L} P\bigl(W_\ell\mid W_{<\ell},X,Z\bigr).
  \label{eq:block_fact}
\end{equation}
In every dataset we consider, the continuous mediators form a single block
placed last, so~\eqref{eq:block_fact} takes the form
\begin{equation}
  P(W\mid X,Z)=
  \underbrace{\prod_{\ell\in\mathcal L_{\mathrm{disc}}}
     P\bigl(W_\ell\mid W_{<\ell},X,Z\bigr)}_{P(W_{\mathrm{disc}}\mid X,Z)}
  \;P\bigl(W_{\mathrm{cont}}\mid W_{\mathrm{disc}},X,Z\bigr).
  \label{eq:cond_fact}
\end{equation}
For ACS, for example, the blocks are
$[\textsc{Educ}]\to[\textsc{Occ}]\to[\textsc{Hours},\textsc{Weeks}]$
(Figure~\ref{fig:block_model}). Placing the discrete blocks upstream of the
continuous one is a substantive ordering claim, made per dataset in
Appendix~\ref{app:roles}, not a convention: it matters for interventions, as
discussed below.

\paragraph{Discrete blocks.}
Each discrete block is modelled by a neural categorical head $g_\phi$ whose
class probabilities depend on $(x,z)$ and on all earlier blocks. Within a block
with several discrete coordinates, the heads are factorized
autoregressively. This is a statistical parameterization of the joint block
distribution and adds no causal edges.

\paragraph{Continuous block.}
The continuous block is modelled by a conditional normalizing flow $f_\theta$
with rational-quadratic spline transforms~\citep{durkan2019neural}, which
captures the joint dependence among continuous mediators without imposing an
order on them. Let $h=[\mathrm{emb}(w_{\mathrm{disc}});\,x;\,\tilde z]$ be the
conditioning vector, where $\mathrm{emb}(\cdot)$ embeds the discrete mediators
and $\tilde z$ is $z$ standardized with training-split statistics. With
$w_{\mathrm{cont}}=f_\theta(u;h)$ and $u\sim\mathcal N(0,I)$,
\begin{equation}
  \log p_\theta(w_{\mathrm{cont}}\mid w_{\mathrm{disc}},x,z)
  =\log\mathcal N\!\bigl(f_\theta^{-1}(w_{\mathrm{cont}};h);0,I\bigr)
  +\log\Bigl|\det\frac{\partial f_\theta^{-1}}{\partial w_{\mathrm{cont}}}\Bigr|.
  \label{eq:flow_loglik}
\end{equation}
The coordinate order inside the flow is a density parameterization and carries
no causal meaning. When a dataset's graph excludes a direct effect of $X$ on
the continuous block, we drop $x$ from $h$. Both factors are trained jointly by
maximizing
\begin{equation}
  \max_{\phi,\theta}\;\frac1n\sum_{i=1}^n\Bigl[
    \sum_{\ell\in\mathcal L_{\mathrm{disc}}}
      \log p_\phi\bigl(w_{\ell,i}\mid w_{<\ell,i},x_i,z_i\bigr)
    +\log p_\theta\bigl(w_{\mathrm{cont},i}\mid w_{\mathrm{disc},i},x_i,z_i\bigr)
  \Bigr].
  \label{eq:flow_training}
\end{equation}
Because every downstream quantity is a Monte Carlo average over model samples,
we select checkpoints by posterior-predictive fidelity on the validation split
rather than by likelihood alone (Appendix~\ref{app:training}).

\paragraph{Monte Carlo estimation.}
Sampling is ancestral: discrete blocks are drawn in order from $g_\phi$, and
$w_{\mathrm{cont}}=f_\theta(u;h)$ is then computed from a fresh
$u\sim\mathcal N(0,I)$. Each nested counterfactual prediction is estimated by
sampling directly from the mediator law of its mediator arm,
\begin{equation}
  \hat\mu_{ab}(z)=\frac1K\sum_{k=1}^{K}\hat f\bigl(a,w^{(k)}_b,z\bigr),
  \qquad w^{(k)}_b\sim\hat P(W\mid X=b,z),
  \label{eq:mu_mc}
\end{equation}
and $\widehat\NDE$ and $\widehat\NIE$ follow by plugging
$\hat\mu_{ab}$ into~\eqref{eq:effects}. We draw separate samples for each
$b$ rather than reweighting samples from one group with importance ratios
$\hat p(w\mid x^{(1)},z)/\hat p(w\mid x^{(0)},z)$, whose variance is large when
the two group-conditional laws overlap poorly. We use
$K=500$ Monte Carlo samples per instance and report percentile-bootstrap
intervals over instances with $B=2000$ replicates. These intervals hold the
fitted mediator model and classifier fixed. Full-pipeline uncertainty is
assessed separately across seeds (Section~\ref{sec:experiments}).

\paragraph{The block order matters only for recourse.}
By the chain rule, any block order represents the same joint distribution
$P(W\mid X,Z)$. The estimands $\mu_{ab}(z)$, and the advantaged reference
$R_z$ used below, therefore do not depend on the order, up to estimation
error, and need only Assumption~\ref{ass:ident}. The order matters only when
we \emph{intervene} on a mediator and regenerate its consequences, which
requires the following stronger assumption.

\begin{assumption}[Interventional validity of the block model]
\label{ass:intervene}
(i)~The blocks follow a valid topological order: no block causes an earlier
block.
(ii)~Modularity: each mechanism $P(W_\ell\mid W_{<\ell},X,Z)$ is unchanged when
earlier blocks are intervened on.
(iii)~There is no unmeasured confounding between mediator blocks given
$(X,Z)$.
(iv)~Intervened values lie within the observed support of $W$.
\end{assumption}

\noindent Assumption~\ref{ass:intervene} is the strongest assumption in the
framework, and only the recourse step uses it. The pre-recourse disparity
needs Assumption~\ref{ass:ident} alone.

% ---------------------------------------------------------------------------
\subsection{Interventional Recourse Distribution}
\label{sec:intervention}

\paragraph{Recipients and actionability.}
Recourse is offered to disadvantaged-group true negatives: instances
$(x^{(0)},w_0,z)$ with $y=0$ and $\hat f(x^{(0)},w_0,z)<\tau$. The values of $X$
and $Z$ are never changed. Each mediator $j$ is assigned an actionability type
that says in which direction an individual can change it:
\emph{fixed} (cannot be changed), \emph{monotone-up} (can only increase, e.g.,
education or hours worked), \emph{monotone-down} (can only decrease), or
\emph{free} (can move in either direction, e.g., occupation). The type
determines the set $\mathcal G_j(w_{0j})$ of values the mediator may take from
its factual value $w_{0j}$. For a discrete mediator, $\mathcal G_j$ contains
the factual category and, for monotone-up, every higher category in its
natural order, or, for free, every category. For a continuous mediator,
$\mathcal G_j$ is a grid with a fixed step between the factual value and a
plausible bound, in the permitted direction (e.g., weekly hours in steps of
5 up to 60). Every grid contains $w_{0j}$, so leaving a mediator unchanged is
always allowed. Appendix~\ref{app:roles} lists each mediator's type, step,
and bounds. An action plan is a target vector
$a=w'\in\mathcal A(w_0):=\prod_j\mathcal G_j(w_{0j})$. Its acted set is
$S(a)=\{j:w'_j\ne w_{0j}\}$, and $\ell^*(a)$ is the earliest block containing
an acted coordinate.

\paragraph{Interventional distribution.}
An action changes the acted mediators and, through the block mechanisms, the
mediators downstream of them. We represent this as a distribution rather than
a point. Let $Q_a(\cdot\mid w_0,x^{(0)},z)$ be the law of $W$ generated as
follows:
\begin{enumerate}[nosep]
  \item blocks before $\ell^*(a)$ keep their factual values;
  \item acted coordinates are clamped, $W_j=w'_j$ for $j\in S(a)$;
  \item unacted coordinates in block $\ell^*(a)$ keep their factual values;
  \item every block after $\ell^*(a)$ is redrawn in order from the learned
        mechanism $\hat P(W_\ell\mid W_{<\ell},x^{(0)},z)$, and acted coordinates
        in that block keep their clamped values.
\end{enumerate}
If $S(a)=\emptyset$, then $Q_a=\delta_{w_0}$. Step~4 regenerates every later
block, not only the true graph descendants of the acted coordinates. This is
conservative and consistent with Assumption~\ref{ass:intervene}: if a later
block does not depend on the acted one, the learned mechanism should reflect
that, and the only cost is added variance. Step~3 resolves the one semantic
choice that blocks introduce. With no causal order inside a block, acting on
one coordinate gives no reason to change the others, so they keep the
individual's values. (Appendix~\ref{app:training} reports a variant in which
they are redrawn from their conditional marginal.)

$Q_a$ is an individualized \emph{interventional} distribution, not a
unit-level counterfactual. We do not abduct exogenous noise, which is not
identifiable for discrete mediators in any case. This follows the
probabilistic view of recourse under imperfect causal
knowledge~\citep{karimi2020imperfect}, and it is why recourse is evaluated over
a distribution of outcomes rather than at a single point.

\paragraph{Prediction distributions.}
Write $\hat f_0(w,z):=\hat f(x^{(0)},w,z)$. For a plan $a$, let
\begin{equation}
  P_a\;:=\;\mathrm{Law}\bigl(\hat f_0(W,z)\bigr),\quad W\sim Q_a(\cdot\mid w_0,x^{(0)},z),
  \label{eq:plan_law}
\end{equation}
be the distribution of the recipient's predicted score after acting. Let the
\emph{advantaged reference} be
\begin{equation}
  R_z\;:=\;\mathrm{Law}\bigl(\hat f_0(W,z)\bigr),\quad W\sim P(W\mid X=x^{(1)},z),
  \qquad
  \mu_{\mathrm{ref}}(z):=\mathbb E[R_z]=\mu_{x^{(0)}x^{(1)}}(z).
  \label{eq:reference}
\end{equation}
$R_z$ is the score distribution the recipient would receive if their mediators
were distributed as those of advantaged individuals with the same $z$. Both
$P_a$ and $R_z$ hold $X$ at $x^{(0)}$ inside $\hat f$, so they differ only
through the mediators: comparing them isolates the mediated pathway and leaves
the direct effect of $X$ untouched. $R_z$ is a population-level counterpart
conditioned on $z$, whereas $Q_a$ retains the recipient's factual
non-descendants. The comparison is therefore between a partly individual
distribution and a $z$-conditional population distribution.

% ---------------------------------------------------------------------------
\subsection{Distribution-Constrained Recourse}
\label{sec:recourse_selection}

\paragraph{Recourse cost.}
The cost of a plan $a=w'$ is
\begin{equation}
  c(a)=\sum_{j\in W_{\mathrm{cont}}}\alpha_j\,(w'_j-w_{0j})^2
  +\sum_{j\in W_{\mathrm{disc}}}\alpha_j\,\mathbf 1\{w'_j\ne w_{0j}\},
  \qquad
  \alpha_j=\begin{cases}
    1/\hat\sigma_j^2 & j\in W_{\mathrm{cont}},\\
    1/H_j & j\in W_{\mathrm{disc}},
  \end{cases}
  \label{eq:cost}
\end{equation}
where $\hat\sigma^2_j$ is the empirical variance of continuous mediator $j$
and $H_j=-\sum_k\hat p_{jk}\log\hat p_{jk}$ is the empirical entropy of discrete
mediator $j$, both computed on the training split. The continuous term is a
squared change in standard-deviation units. The discrete term makes changes to
widely spread mediators cheaper than changes to concentrated ones.

\paragraph{Constrained program.}
A plan is \emph{$\gamma$-valid} if it yields a positive decision with
probability at least $\gamma$, i.e., $P_a([\tau,1])\ge\gamma$. Our method
solves
\begin{equation}
  a^*\in\arg\min_{a\in\mathcal A(w_0)}\;c(a)
  \quad\text{s.t.}\quad
  \underbrace{P_a([\tau,1])\ge\gamma}_{\text{validity}},
  \qquad
  \underbrace{W_1(P_a,R_z)\le\varepsilon}_{\text{distributional parity}},
  \label{eq:program}
\end{equation}
where $W_1$ is the 1-Wasserstein distance on $[0,1]$. If the feasible set is
empty, the method \emph{abstains}: the recipient receives the zero-cost
factual plan $a=w_0$ and the abstention is recorded. Abstaining recipients
stay in every evaluation average, so difficult cases are never silently
dropped.

Validity alone asks that the recipient \emph{cross} the decision boundary. The
parity constraint asks that the recipient's score distribution \emph{land
where their advantaged counterpart's would be}. It penalizes undershooting
(valid, but still scored below the reference) as well as overshooting (pushed
further than the reference). Placing the Wasserstein constraint in prediction
space has two advantages. First, $\hat f_0$ maps both mediator distributions to
$[0,1]$, so we need no metric over mixed continuous and categorical
mediators, and no Lipschitz assumption on $\hat f$ over categorical codes.
Second, in one dimension $W_1$ has the closed form
$W_1(G,R)=\int_0^1|G^{-1}(u)-R^{-1}(u)|\,du$, which yields the exact analysis
in Section~\ref{sec:theory}.

\paragraph{Baseline.}
The baseline, \emph{ordinary actionable recourse}, solves~\eqref{eq:program}
without the parity constraint. We deliberately make it as strong as we can,
so that the comparison is the hardest one for our method rather than a
comparison against a straw man. Actionable recourse as usually
formulated~\citep{ustun2019actionable} changes features directly, treats them
as independently manipulable, and checks validity at a single point. Our
baseline instead inherits every component of our method except the parity
constraint:
(i)~the same action space, restricted to the mediators with the same
actionability types and grids;
(ii)~the same causal intervention model $Q_a$, so downstream consequences of
an action are propagated through the learned block mechanisms rather than
ignored~\citep{karimi2021algorithmic};
(iii)~the same probabilistic notion of validity, $P_a([\tau,1])\ge\gamma$,
evaluated over the same interventional draws;
(iv)~the same cost~\eqref{eq:cost}; and
(v)~the same abstention rule.
Both methods also select from identical candidate and reference draws, so
neither benefits from Monte Carlo noise. Any difference between them is
therefore due to the parity constraint alone. In
Section~\ref{sec:exp_gamma} we go one step further and give the baseline its
best validity level, chosen after the fact on the evaluation recipients,
while our method keeps a fixed $\gamma$.

\paragraph{Estimation.}
For each recipient we enumerate $\mathcal A(w_0)$ in full. For each plan we draw
$K_{\mathrm{int}}$ samples from $Q_a$ and estimate validity as the fraction of
scores at or above $\tau$. We also draw $K_{\mathrm{ref}}$ samples from
$\hat P(W\mid x^{(1)},z)$ to form the empirical reference, and we evaluate
$W_1$ between the two empirical score distributions. The guarantees of
Section~\ref{sec:theory} hold exactly for these empirical distributions. The
empirical $W_1$ is biased upward at rate $O(K^{-1/2})$, so $\varepsilon$ should
be set above the Monte Carlo noise level. We use $\gamma=0.5$, $\varepsilon=0.1$,
$K_{\mathrm{int}}=32$, and $K_{\mathrm{ref}}=200$, and we vary $\gamma$ in
Section~\ref{sec:exp_gamma}.
Appendix~\ref{app:soft} describes a soft penalized objective (cost plus
weighted shortfall, pointwise direct-effect, and transport terms) that the
ablations in Section~\ref{sec:experiments} compare against.

% ---------------------------------------------------------------------------
\subsection{When Are Validity and Distributional Parity Compatible?}
\label{sec:theory}

The two constraints in~\eqref{eq:program} can conflict: if the advantaged
reference itself rarely clears the threshold, any sufficiently valid plan must
move the recipient \emph{beyond} it. The following result makes this conflict
precise.

\begin{theorem}{Validity--Parity Tradeoff}{val-par}
Fix a recipient with covariates $z$ and factual mediators $w_0$, and define
\[
\varepsilon_\gamma(z)
:=
\inf_{G:\,G([\tau,1])\ge\gamma}
W_1(G,R_z)
\]
as the smallest distributional parity gap attainable by any score distribution satisfying the validity requirement. We define the admissible action set as follows:

\[
\varepsilon_{\mathcal A}(z,w_0)
:=
\min_{\substack{a\in\mathcal A(w_0):\\P_a([\tau,1])\ge\gamma}}
W_1(P_a,R_z),
\]

with $\varepsilon_{\mathcal A}(z,w_0)=+\infty$ if no $\gamma$-valid action exists. Then
\begin{equation}
\varepsilon_{\mathcal A}(z,w_0)
\ge
\varepsilon_\gamma(z).
\label{eq:feasibility_bound}
\end{equation}

Hence, validity imposes an irreducible lower bound on attainable distributional parity. Therefore, the recourse program is feasible only when the following bound holds:

\[
\varepsilon
\ge
\varepsilon_{\mathcal A}(z,w_0)
\ge
\varepsilon_\gamma(z).
\]

When the program yields a feasible action, its solution satisfies the following bound over the remaining parity gap:

\[
W_1(P_{a^\ast},R_z)\le\varepsilon,
\]
 Moreover, $\varepsilon_\gamma(z)=0$ if and only if the reference distribution itself satisfies $R_z([\tau,1])\ge\gamma$.
\end{theorem}

\noindent The proof is in Appendix~\ref{app:theory}, which also shows that
the infimum defining $\varepsilon_\gamma(z)$ is attained and has the closed form
$\varepsilon_\gamma(z)=\int_{1-\gamma}^{1}(\tau-F_{R_z}^{-1}(u))_+\,du$, where
$F_{R_z}^{-1}$ is the quantile function of $R_z$. It can therefore be computed
exactly from samples of $R_z$, and it depends only on $\hat f$, the advantaged
mediator distribution at $z$, and $(\tau,\gamma)$, not on the action set or
the cost.

By~\eqref{eq:feasibility_bound}, every abstention has exactly one of three
causes:
(a)~\emph{no recourse}: no $\gamma$-valid plan exists;
(b)~\emph{target incompatibility}: $\varepsilon<\varepsilon_\gamma(z)$, a
property of $\hat f$, $P(W\mid x^{(1)},z)$, and $\gamma$ alone; or
(c)~\emph{unreachability}: $\varepsilon_\gamma(z)\le\varepsilon<
\varepsilon_{\mathcal A}(z,w_0)$, a property of the action set and the
mechanisms.
Appendix~\ref{app:theory} proves two further consequences. First, every
$\gamma$-valid plan has $\mathbb E[P_a]\ge\gamma\tau$, so when
$\mu_{\mathrm{ref}}(z)<\gamma\tau$ validity alone forces every plan, including
ordinary recourse, to overshoot the reference
(Proposition~\ref{prop:overshoot}). Second, since feasible recipients satisfy
$|D_{\mathrm{post},i}|\le\varepsilon$ and abstaining recipients keep their
factual disparity, closure is limited by the abstention rate
(Proposition~\ref{prop:guarantee}), and (a)--(c) explain what drives
abstention.

% ---------------------------------------------------------------------------
\subsection{Estimands and Evaluation}
\label{sec:evaluation}

\paragraph{Pre-recourse decomposition.}
On the classifier scale, using~\eqref{eq:mu_ab}, we report the pure natural
direct and indirect effects,
\begin{equation}
  \NDE(z)=\mu_{10}(z)-\mu_{00}(z),\qquad
  \NIE(z)=\mu_{01}(z)-\mu_{00}(z),
  \label{eq:effects}
\end{equation}
where indices $0,1$ abbreviate $x^{(0)},x^{(1)}$, estimated
via~\eqref{eq:mu_mc} and averaged over held-out instances.

Both $P_a$ and $R_z$ fix $X=x^{(0)}$ inside $\hat f$, so recourse on $W$ can at
most close the mediated part of the disparity, the $\NIE$. The $\NDE$ is
carried along $X\to\hat Y$ and no mediator-based recourse can reach it.

\paragraph{Mediated disparity before and after recourse.}
For recipient $i$ with confounders $z_i$, we define
\begin{align}
  D_{\mathrm{pre},i}&:=\mu_{\mathrm{ref}}(z_i)-\mu_{00}(z_i)=\NIE(z_i),
  \label{eq:dpre}\\
  D^{\mathrm{fact}}_{\mathrm{pre},i}&:=\mu_{\mathrm{ref}}(z_i)-\hat f_0(w_{0,i},z_i),
  \label{eq:dpre_fact}\\
  D_{\mathrm{post},i}&:=\mu_{\mathrm{ref}}(z_i)-\mathbb E\bigl[P_{a^*_i}\bigr].
  \label{eq:dpost}
\end{align}
$D_{\mathrm{pre}}$ compares the reference with the disadvantaged group's natural
mediator distribution. $D^{\mathrm{fact}}_{\mathrm{pre}}$ compares it with the
recipient's own factual score, the natural baseline for an individual
intervention. $D_{\mathrm{post}}$ compares a natural distribution with an
\emph{intervened} one, so it is not a natural indirect effect, and we do not
call it one. The reference is never itself passed through recourse.

\paragraph{Metrics.}
Our primary metric is recipient-level closure,
\begin{equation}
  \mathrm{Closure}\;=\;1-\frac{\sum_i|D_{\mathrm{post},i}|}
                             {\sum_i|D^{\mathrm{fact}}_{\mathrm{pre},i}|},
  \label{eq:closure}
\end{equation}
computed from absolute values per recipient so that positive and negative
gaps cannot cancel. We also report:
the overshoot rate, $\Pr[D_{\mathrm{post},i}<0]$, and the mean overshoot size;
the paired improvement over the baseline,
$|D^{\mathrm{base}}_{\mathrm{post},i}|-|D_{\mathrm{post},i}|$;
the abstention rate;
and the mean recourse cost among feasible recipients.

\paragraph{Compatibility diagnostics.}
To tell \emph{no recourse exists} apart from \emph{the targets conflict}, we
report, per dataset and classifier:
the share of recipients whose reference is itself $\gamma$-valid, i.e.,
$\varepsilon_\gamma(z)=0$;
the distribution of $\varepsilon_\gamma(z)$ and of
$\mu_{\mathrm{ref}}(z)-\gamma\tau$;
the distribution of $\varepsilon_{\mathcal A}(z,w_0)$, whose empirical CDF gives
coverage as a function of $\varepsilon$;
and the breakdown of abstentions into causes~(a)--(c).
All four are computed from the same empirical score distributions used to
solve~\eqref{eq:program}.
